\documentclass[11pt,a4paper]{article}

%----------------------------------------------------------------------------------------
%	PACKAGES
%----------------------------------------------------------------------------------------
\usepackage[top=1.8cm, bottom=1.8cm, left=2cm, right=2cm]{geometry}
\usepackage[utf8]{inputenc}
\usepackage[T1]{fontenc}
\usepackage[default]{lato}
\usepackage{xcolor}
\usepackage{fontawesome5}
\usepackage{titlesec}
\usepackage{enumitem}
\usepackage{hyperref}
\usepackage{tabularx}

%----------------------------------------------------------------------------------------
%	COLORS & STYLES
%----------------------------------------------------------------------------------------
\definecolor{accentred}{HTML}{E03C31}
\definecolor{darktext}{HTML}{1A1A1A}
\definecolor{graytext}{HTML}{666666}

\pagestyle{empty}

% Hyperref setup
\hypersetup{
    colorlinks=true,
    linkcolor=accentred,
    urlcolor=accentred,
    pdftitle={Eden Gilbert Kiseka - Resume},
    pdfauthor={Eden Gilbert Kiseka}
}

% Section Styling
\titleformat{\section}
  {\vspace{1em}\large\bfseries\color{darktext}} % format
  {} % label
  {0em} % sep
  {} % before
  [\vspace{-0.4em}\rule{\textwidth}{0.4pt}] % after (line)

% Itemize settings
\setlist[itemize]{leftmargin=1.5em, nosep, topsep=0.2em, itemsep=0.2em}

% Job Entry Command
\newcommand{\jobentry}[5]{
    \noindent
    \begin{minipage}[t]{0.75\textwidth}
        \textbf{\color{accentred} #1}
    \end{minipage}%
    \begin{minipage}[t]{0.25\textwidth}
        \raggedleft
        \color{graytext}\small #2
    \end{minipage}\\[0.1em]
    \noindent {\color{graytext} #3} \hfill {\color{graytext}\footnotesize #4} \\[0.3em]
    \noindent \small #5
    \vspace{1.2em}
}

% Project Entry Command
\newcommand{\projectentry}[4]{
    \noindent
    \begin{minipage}[t]{0.75\textwidth}
        \textbf{\color{accentred} #1}
    \end{minipage}%
    \begin{minipage}[t]{0.25\textwidth}
        \raggedleft
        \color{graytext}\small \href{#4}{View Repository}
    \end{minipage}\\[0.1em]
    \noindent {\color{graytext}\small #2} \\[0.3em]
    \noindent \small #3
    \vspace{1.2em}
}

% Skill Entry Command
\newcommand{\skillentry}[2]{
    \noindent \textbf{#1} \hfill \color{graytext} #2 \\
}

\begin{document}

%--- HEADER ---
\noindent
\begin{minipage}[t]{0.6\textwidth}
    {\Huge \bfseries \color{darktext} Eden Gilbert Kiseka} \\[0.4em]
    {\Large \bfseries \color{accentred} Embedded Systems Developer \& Technical Educator}
\end{minipage}%
\begin{minipage}[t]{0.4\textwidth}
    \raggedleft \small \color{darktext}
    \faPhone \hspace{0.5em} +256 703 055 006 \\
    \faEnvelope \hspace{0.5em} \href{mailto:edengilbertus@proton.me}{edengilbertus@proton.me} \\
    \faGlobe \hspace{0.5em} \href{https://edengilbert.is-a.dev}{edengilbert.is-a.dev} \\
    \faLinkedin \hspace{0.5em} \href{https://linkedin.com/in/edengilbert}{edengilbert} \\
    \faGithub \hspace{0.5em} \href{https://github.com/edengilbertus}{edengilbertus} \\
    \faMapMarker \hspace{0.5em} Kampala, Uganda
\end{minipage}

\vspace{1.5em}

%--- EXPERIENCE ---
\section*{EXPERIENCE}

\jobentry{Backend Engineer}{Oct 2025 -- Present}
{Googah Goats Limited}{Kampala, Uganda}
{
    \begin{itemize}
        \item Architect backend infrastructure using Google AppScript and JavaScript.
        \item Design automation systems and data workflows managing business operations.
        \item Build custom APIs and integration scripts connecting various services.
        \item Develop understanding of distributed systems and API design.
    \end{itemize}
}

\jobentry{Technical Instructor -- Embedded Systems}{Oct 2024 -- Sep 2025}
{Neriko Electronics}{Kampala, Uganda}
{
    \begin{itemize}
        \item Taught embedded systems programming to students using C, C++, and Arduino platform.
        \item Developed comprehensive curriculum covering microcontroller fundamentals through advanced IoT applications.
        \item Designed hands-on laboratory exercises demonstrating real-world embedded applications.
        \item Covered topics: GPIO programming, UART/SPI/I2C communication, ADC/DAC interfacing, PWM control, interrupt handling, timer programming.
        \item Instructed students in industry tools: PlatformIO (embedded development), TinkerCAD (circuit simulation), Fritzing (PCB design), Proteus (circuit simulation), AutoCAD Electrical (professional schematics).
        \item Created educational materials: lecture notes, circuit diagrams, code examples, project specifications.
        \item Mentored students individually through debugging hardware-software integration challenges.
        \item Taught systematic troubleshooting using multimeters, oscilloscopes, and logic analyzers.
    \end{itemize}
}

%--- PROJECTS ---
\section*{PROJECTS}

\projectentry{kore -- Schr\"{o}dinger-Poisson Solver}{Rust, PyO3, C ABI, NumPy, RustFFT, Rayon}{
    \begin{itemize}
        \item Portable pseudo-spectral kernel for fuzzy dark matter simulation workloads.
        \item Implements split-step Fourier evolution with RustFFT and Rayon threading for parallelism.
        \item Exposes a stable C ABI and \texttt{pykore} Python bindings via PyO3/NumPy, bridging Python workflows with HPC-grade performance.
        \item Verified with Rust integration tests, strict C11 smoke builds, and Python tests covering in-place complex128 stepping with GIL release on the compute path.
    \end{itemize}
}{https://github.com/edengilbertus/kore}

\projectentry{Amber -- Android Reconstruction Pipeline}{Kotlin, DEX IR, SSA, JavaPoet, Gradle}{
    \begin{itemize}
        \item Lifts APK and DEX artifacts into typed IR, CFG, SSA, and generated source, then synthesizes a complete Android Studio project around the result.
        \item Combines dex parsing, Java/Kotlin emission, Gradle synthesis, metadata analysis, and evidence-rich diagnostics to turn reverse engineering into a structured engineering workflow.
        \item Supports end-to-end reconstruction: emitted Java baselines, Gradle project generation, metadata-aware Kotlin recovery, Compose detection, and resource pipeline integration.
        \item Open source under Apache 2.0; CI-driven validation on real APKs.
    \end{itemize}
}{https://github.com/edengilbertus/amber}

%--- SKILLS ---
\section*{TECHNICAL SKILLS}

\begin{minipage}[t]{0.48\textwidth}
    \skillentry{Embedded Programming}{C, C++, Arduino}
    \skillentry{Microcontroller Tools}{Microchip Studio, PlatformIO, Arduino IDE}
    \skillentry{Communication Protocols}{UART, SPI, I2C, GPIO, PWM}
    \skillentry{Peripherals}{ADC/DAC, Timers, Interrupts, Sensors}
    \skillentry{Circuit Design}{TinkerCAD, Fritzing, Proteus, AutoCAD Electrical}
    \skillentry{Debug Tools}{Multimeter, Oscilloscope, Logic Analyzer}
    \skillentry{Platforms}{Arduino (Uno, Mega, Nano), ESP32/ESP8266, STM32}
    \skillentry{Embedded Linux}{Raspberry Pi, Cross-Compilation, Device Drivers}
\end{minipage}%
\hfill
\begin{minipage}[t]{0.48\textwidth}
    \skillentry{Software Development}{Python, Kotlin, Java, JavaScript, Rust, Swift}
    \skillentry{Digital Design}{Logisim, FPGA (learning), RISC-V, Logic Simulation}
    \skillentry{Hardware Description}{Verilog/SystemVerilog (learning)}
    \skillentry{Version Control}{Git, GitHub, GitLab}
    \skillentry{Development IDEs}{VS Code, PlatformIO, Microchip Studio, Arduino IDE}
    \skillentry{Operating Systems}{Linux, Windows}
    \skillentry{Build Systems}{Makefiles, CMake, PlatformIO}
\end{minipage}

\vspace{1em}

%--- EDUCATION ---
\section*{EDUCATION}

\jobentry{Associate Degree in Electrical Engineering \& Computer Science}{2025 -- 2027}
{University of the People}{United States}
{}

\jobentry{Diploma in Electronics and Electrical Engineering}{2025 -- 2027}
{Uganda Institute of ICT}{Uganda}
{}

%--- LANGUAGES ---
\section*{LANGUAGES}

\noindent
\begin{minipage}[t]{0.32\textwidth}
    \textbf{English} \hfill {\color{graytext} Fluent}
\end{minipage}%
\hfill
\begin{minipage}[t]{0.32\textwidth}
    \textbf{Luganda} \hfill {\color{graytext} Fluent}
\end{minipage}%
\hfill
\begin{minipage}[t]{0.32\textwidth}
    \textbf{French} \hfill {\color{graytext} Intermediate (B2)}
\end{minipage}

\end{document}
